\documentclass[uplatex,dvipdfmx,a4paper,11pt]{jsarticle}
\input{../../../templates/note.tex}
\noteprint{数理物理1　第8章「行列の対角化と基準振動」}{2026-11-25}
%% note-moto	problems\数理物理1\shou08\数理物理1-8.tex
\begin{document}
\begin{kirokuichiran}
  \ichirangyou{\hyperlink{bt-m-054-orth3}{1}}{L1・計算}{◎◎◎}{S}{2026-10-03}
  \ichirangyou{\hyperlink{bt-m-055-unitary2}{2}}{L1・計算}{◎◎◎}{S}{2026-10-03}
  \ichirangyou{\hyperlink{bt-m-056-hermite2}{3}}{L1・計算}{×◎◎}{C}{2026-10-03 (1)未着手}
  \ichirangyou{\hyperlink{bt-m-061-ring3}{4}}{L2・計算}{◎△××}{D}{2026-10-03 (2)計算 (3)未着手 (4)未着手}
  \ichirangyou{\hyperlink{bt-m-063-basis2}{5}}{L2・計算}{◎◎◎◎}{S}{2026-10-03}
  \ichirangyou{6}{L3・計算}{}{未着手}{}
  \ichirangyou{7}{L3・導出}{}{未着手}{}
  \ichirangyou{8}{L4・計算}{}{未着手}{}
\end{kirokuichiran}

\youtenseiri
\begin{noteshikilist}
  \item Dirac の記法の縦ベクトルと、その $\dagger$（複素共役をとった横ベクトル） % 出典: p.66--67（間の式 7.1.76, 7.1.80--7.1.81, 7.1.84 は載せない）
\begin{align*}
&|f\rangle,\ |g\rangle \notetag{8.1}\\
&\langle f|,\ \langle g| \notetag{8.2}
\end{align*}
  \item 内積の複素共役、基底ベクトルの正規直交性と完備性 % 出典: p.67
\begin{align*}
\langle f|g\rangle^{*}&=\langle g|f\rangle \notetag{8.3}\\
\langle n|m\rangle&=\delta_{nm} \notetag{8.4}\\
\sum_{n=1}^{N}|n\rangle\langle n|&=1 \notetag{8.5}
\end{align*}
  \item 別の基底 $|\mu\rangle$ による展開・成分と変換行列 % 出典: p.67
\begin{align*}
|f\rangle&=\sum_{\mu=1}^{N}|\mu\rangle\langle\mu|f\rangle \notetag{8.6}\\
\hat{f}_{\mu}&=\langle\mu|f\rangle \notetag{8.7}\\
f_n&=\sum_{\mu=1}^{N}\langle n|\mu\rangle\hat{f}_{\mu} \notetag{8.8}\\
U_{n\mu}&=\langle n|\mu\rangle,\qquad (U^{\dagger})_{\mu n}=\langle\mu|n\rangle \notetag{8.9}
\end{align*}
  \item $U$ のユニタリー性 % 出典: p.68
\begin{align*}
UU^{\dagger}&=U^{\dagger}U=I \notetag{8.10}\\
(UU^{\dagger})_{nm}&=\sum_{\mu=1}^{N}U_{n\mu}(U^{\dagger})_{\mu m}=\sum_{\mu=1}^{N}\langle n|\mu\rangle\langle\mu|m\rangle=\delta_{nm} \notetag{8.11}\\
(U^{\dagger}U)_{\mu\nu}&=\sum_{n=1}^{N}(U^{\dagger})_{\mu n}U_{n\nu}=\sum_{n=1}^{N}\langle\mu|n\rangle\langle n|\nu\rangle=\delta_{\mu\nu} \notetag{8.12}
\end{align*}
\end{noteshikilist}

\mondaihen
\level{1}
% toubun: kiroku\toubun\数理物理1\054_2026-10-03.tex
\mondaibangou{1}
\begin{mondai}{054-orth3}{計算}{3次の直交行列による直交変換}
行列
\[
C=\frac{1}{3}\begin{pmatrix}1&2&2\\2&1&-2\\2&-2&1\end{pmatrix}
\]
を考える。
\begin{shoumon}
  \item $C^{t}C=CC^{t}=I$ を示せ。
  \item $\vect{f}=(3,3,3)^{t}$ に対して $\hat{\vect{f}}=C^{t}\vect{f}$ を求め、$\abs{\vect{f}}^{2}=\abs{\hat{\vect{f}}}^{2}$ を確かめよ。
  \item $C$ の第 $\mu$ 列を列ベクトル $\vect{C}_{\mu}$（$\mu=1,2,3$）とする。次の式が成り立つことを成分で確かめよ。
\[
\vect{f}=\sum_{\mu=1}^{3}\vect{C}_{\mu}\hat{f}_{\mu}
\]
\end{shoumon}
\notekiroku{2026-10-03 S ◎◎◎}
\begin{kotaeutsushi}{2026-10-03}
\toi{1}{◎}{}
\begin{align*}
  C^{t}C&=\frac{1}{9}\begin{pmatrix}1&2&2\\2&1&-2\\2&-2&1\end{pmatrix}\begin{pmatrix}1&2&2\\2&1&-2\\2&-2&1\end{pmatrix}\\
  &=\frac{1}{9}\begin{pmatrix}9&0&0\\0&9&0\\0&0&9\end{pmatrix}=I
\end{align*}
\toi{2}{◎}{}
\[
  \hat{\vect{f}}=\kakimachigai{\begin{pmatrix}1&2&2\\2&1&-2\\2&-2&1\end{pmatrix}}{\tfrac{1}{3}(\cdots)}\begin{pmatrix}3\\3\\3\end{pmatrix}=\begin{pmatrix}5\\1\\1\end{pmatrix}
\]
\[
  \abs{\vect{f}}^{2}=(3\ 3\ 3)\begin{pmatrix}3\\3\\3\end{pmatrix}=27,\qquad
  \abs{\hat{\vect{f}}}^{2}=(5\ 1\ 1)\begin{pmatrix}5\\1\\1\end{pmatrix}=27
\]
\toi{3}{◎}{}
\begin{align*}
  &\frac{1}{3}\biggl(5\begin{pmatrix}1\\2\\2\end{pmatrix}+1\begin{pmatrix}2\\1\\-2\end{pmatrix}\kakimachigai{\,}{+}1\begin{pmatrix}2\\-2\\1\end{pmatrix}\biggr)\\
  &\qquad=\kakimachigai{\begin{pmatrix}9\\9\\9\end{pmatrix}}{\tfrac{1}{3}(9,9,9)^{t}}=(3,3,3)^{t}
\end{align*}
（脇に $\frac{1}{3}(\fumei,1,1)^{t}$ のメモ）
\end{kotaeutsushi}
\end{mondai}

% toubun: kiroku\toubun\数理物理1\055_2026-10-03.tex
\mondaibangou{2}
\begin{mondai}{055-unitary2}{計算}{2次のユニタリー行列}
行列
\[
U=\frac{1}{\sqrt{2}}\begin{pmatrix}1&\ii\\ \ii&1\end{pmatrix}
\]
を考える。
\begin{shoumon}
  \item $U^{\dagger}$ を求め、$U^{\dagger}U=UU^{\dagger}=I$ を示せ。
  \item $\vect{f}=(1,\ii)^{t}$ に対して $\hat{\vect{f}}=U^{\dagger}\vect{f}$ を求め、$\abs{\vect{f}}^{2}=\abs{\hat{\vect{f}}}^{2}$ を確かめよ。
  \item $U^{t}U$ を計算し、$U$ が直交行列ではないことを確かめよ。
\end{shoumon}
\notekiroku{2026-10-03 S ◎◎◎}
\begin{kotaeutsushi}{2026-10-03}
\toi{1}{◎}{}
\[
  U^{\dagger}=\frac{1}{\sqrt{2}}\begin{pmatrix}1&-\ii\\-\ii&1\end{pmatrix}
\]
\begin{align*}
  U^{\dagger}U&=\frac{1}{2}\begin{pmatrix}1&-\ii\\-\ii&1\end{pmatrix}\begin{pmatrix}1&\ii\\\ii&1\end{pmatrix}\\
  &=\frac{1}{2}\begin{pmatrix}2&0\\0&2\end{pmatrix}=I
\end{align*}
\toi{2}{◎}{}
\[
  \vect{f}=\begin{pmatrix}1\\\ii\end{pmatrix}
\]
\begin{align*}
  \hat{\vect{f}}&=\frac{1}{\sqrt{2}}\begin{pmatrix}1&-\ii\\-\ii&1\end{pmatrix}\begin{pmatrix}1\\\ii\end{pmatrix}\\
  &=\frac{1}{\sqrt{2}}\begin{pmatrix}2\\0\end{pmatrix}=\begin{pmatrix}\sqrt{2}\\0\end{pmatrix}
\end{align*}
\[
  \abs{\vect{f}}^{2}=2,\qquad \abs{\hat{\vect{f}}}^{2}=2
\]
\toi{3}{◎}{}
\begin{align*}
  U^{t}U&=\frac{1}{2}\begin{pmatrix}1&\ii\\\ii&1\end{pmatrix}\begin{pmatrix}1&\ii\\\ii&1\end{pmatrix}\\
  &=\frac{1}{2}\begin{pmatrix}0&2\ii\\2\ii&0\end{pmatrix}=\begin{pmatrix}0&\ii\\\ii&0\end{pmatrix}
\end{align*}
\end{kotaeutsushi}
\end{mondai}

% toubun: kiroku\toubun\数理物理1\056_2026-10-03.tex
\mondaibangou{3}
\begin{mondai}{056-hermite2}{計算}{2次のエルミート行列の対角化}
行列
\[
H=\begin{pmatrix}1&2\ii\\-2\ii&1\end{pmatrix}
\]
を考える。
\begin{shoumon}
  \item $H=H^{\dagger}$ を確かめよ。
  \item $H$ の固有値と、規格化された固有ベクトルを求めよ。
  \item (2) の固有ベクトルを列として並べた行列 $A$ がユニタリー行列であることを示し、$A^{\dagger}HA$ を計算して対角行列になることを確かめよ。
\end{shoumon}
\notekiroku{2026-10-03 C ×◎◎}
\begin{kotaeutsushi}{2026-10-03}
\toi{1}{×}{未着手}
\mikai
\toi{2}{◎}{}
\begin{align*}
  \begin{vmatrix}1-\lambda&2\ii\\-2\ii&1-\lambda\end{vmatrix}&=(1-\lambda)^{2}-4\\
  &=\lambda^{2}-2\lambda-3=(\lambda-3)(\lambda+1)
\end{align*}
\[ \lambda=3,\ -1 \]
$\lambda=3$：
\[ \begin{pmatrix}-2&2\ii\\-2\ii&-2\end{pmatrix}\begin{pmatrix}1\\-\ii\end{pmatrix} \]
$\lambda=-1$：
\[ \begin{pmatrix}2&2\ii\\-2\ii&2\end{pmatrix}\begin{pmatrix}1\\\ii\end{pmatrix} \]
\[
  \frac{1}{\sqrt{2}}\begin{pmatrix}1&1\\-\ii&\ii\end{pmatrix},\qquad \begin{pmatrix}3&0\\0&-1\end{pmatrix}
\]
\toi{3}{◎}{}
\[
  A=\frac{1}{\sqrt{2}}\begin{pmatrix}1&1\\-\ii&\ii\end{pmatrix}
\]
\begin{align*}
  A^{\dagger}A&=\frac{1}{2}\begin{pmatrix}1&\ii\\1&-\ii\end{pmatrix}\begin{pmatrix}1&1\\-\ii&\ii\end{pmatrix}\\
  &=\frac{1}{2}\begin{pmatrix}2&0\\0&2\end{pmatrix}=I
\end{align*}
\begin{align*}
  A^{\dagger}HA&=\frac{1}{2}\begin{pmatrix}1&\ii\\1&-\ii\end{pmatrix}\begin{pmatrix}1&2\ii\\-2\ii&1\end{pmatrix}\begin{pmatrix}1&1\\-\ii&\ii\end{pmatrix}\\
  &=\frac{1}{2}\begin{pmatrix}3&3\ii\\-1&\ii\end{pmatrix}\begin{pmatrix}1&1\\-\ii&\ii\end{pmatrix}\\
  &=\frac{1}{2}\begin{pmatrix}6&0\\0&-2\end{pmatrix}=\begin{pmatrix}3&0\\0&-1\end{pmatrix}
\end{align*}
\end{kotaeutsushi}
\begin{kizuki}
\item (1) の $H=H^{\dagger}$ は，転置して複素共役をとった $H^{\dagger}$ を書き，元の $H$ と成分ごとに比べれば終わる。すぐ済む小問ほど飛ばさずに書く。
\item (2) の答えは規格化した形 $\frac{1}{\sqrt{2}}(1,\mp\ii)^{t}$ で書く。規格化の因子を (3) の $A$ でだけ付けると，(2) の答えとしては足りない。
\end{kizuki}
\begin{tadashii}
\houshin エルミート行列 $H$ の固有値・規格化固有ベクトルを求め、固有ベクトルを並べた行列 $A$ で $H$ が対角化されることを確かめる。使う定理：エルミート行列の固有値は実数で、異なる固有値に属する固有ベクトルは直交する。規格化は
\[
\sum_n\abs{A_{n\mu}}^{2}=1
\]
である。固有ベクトルを並べた $A$ はユニタリー行列で、$A^{\dagger}HA$ は固有値を並べた対角行列になる。
\keisan
\shoumonkai{1} 転置して複素共役をとると、$(1,2)$ 成分は $(-2\ii)^{*}=2\ii$、$(2,1)$ 成分は $(2\ii)^{*}=-2\ii$、対角成分は実数 $1$ のままなので $H^{\dagger}=H$。
\shoumonkai{2} 固有方程式は
\[
\det(H-\lambda I)=(1-\lambda)^{2}-(2\ii)(-2\ii)=(1-\lambda)^{2}-4=0
\]
より $\lambda=3,\,-1$（どちらも実数）。
$\lambda=3$ のとき $(1-3)a+2\ii b=0$ より $b=-\ii a$。規格化して $\vect{v}_1=\frac{1}{\sqrt{2}}(1,-\ii)^{t}$。
$\lambda=-1$ のとき $(1+1)a+2\ii b=0$ より $b=\ii a$。規格化して $\vect{v}_2=\frac{1}{\sqrt{2}}(1,\ii)^{t}$。
直交性の確認: $\vect{v}_1^{\dagger}\vect{v}_2=\frac{1}{2}\bigl(1\cdot1+\ii\cdot\ii\bigr)=0$。
\shoumonkai{3}
\[
A=\frac{1}{\sqrt{2}}\begin{pmatrix}1&1\\-\ii&\ii\end{pmatrix},\qquad A^{\dagger}=\frac{1}{\sqrt{2}}\begin{pmatrix}1&\ii\\1&-\ii\end{pmatrix}.
\]
\[
A^{\dagger}A=\frac{1}{2}\begin{pmatrix}1+\ii(-\ii)&1+\ii\cdot\ii\\1+(-\ii)(-\ii)&1+(-\ii)\ii\end{pmatrix}=\frac{1}{2}\begin{pmatrix}2&0\\0&2\end{pmatrix}=I
\]
（正方行列なので $AA^{\dagger}=I$ も成り立つ）。
$HA$ の各列は $3\vect{v}_1$、$-\vect{v}_2$ なので $HA=\frac{1}{\sqrt{2}}\begin{pmatrix}3&-1\\-3\ii&-\ii\end{pmatrix}$。よって
\[
A^{\dagger}HA=\frac{1}{2}\begin{pmatrix}3+\ii(-3\ii)&-1+\ii(-\ii)\\3+(-\ii)(-3\ii)&-1+(-\ii)(-\ii)\end{pmatrix}=\frac{1}{2}\begin{pmatrix}6&0\\0&-2\end{pmatrix}=\begin{pmatrix}3&0\\0&-1\end{pmatrix}.
\]
\kotae{\shoumonkai{1} $H^{\dagger}=H$。\shoumonkai{2} 固有値 $3$: $\vect{v}_1=\frac{1}{\sqrt{2}}(1,-\ii)^{t}$、固有値 $-1$: $\vect{v}_2=\frac{1}{\sqrt{2}}(1,\ii)^{t}$（位相の任意性を除く）。\shoumonkai{3} $A^{\dagger}A=I$、$A^{\dagger}HA=\begin{pmatrix}3&0\\0&-1\end{pmatrix}$。}
\kenzan 固有値 $3,-1$ がどちらも実数であること、$\vect{v}_1^{\dagger}\vect{v}_2=0$ で固有ベクトルが直交すること、$A^{\dagger}A=I$ となること、$A^{\dagger}HA$ の対角成分に固有値が並ぶことを確かめた。これらはエルミート行列の性質と整合する。
\tsumazuki 固有ベクトルには絶対値1の定数倍（位相）の任意性がある。
\end{tadashii}
\end{mondai}

\level{2}
% toubun: kiroku\toubun\数理物理1\061_2026-10-03.tex
\mondaibangou{4}
\begin{mondai}{061-ring3}{計算}{3質点の環状ばね系の基準振動}
質量 $m$ の3個の質点が円周上に並び、隣り合う質点どうしがばね定数 $k$ のばねでつながれている（質点1と2、2と3、3と1の3本）。さらに各質点は、ばね定数 $\kappa$（$\kappa>0$）のばねで自分の釣り合いの位置にもつながれている。質点は円周に沿った1方向にだけ動き、質点 $n$ の釣り合いの位置からの変位を $x_n$（$n=1,2,3$）とする。ポテンシャルエネルギーは
\[
V=\frac{\kappa}{2}\sum_{n=1}^{3}x_n^{2}+\frac{k}{2}\bigl\{(x_2-x_1)^{2}+(x_3-x_2)^{2}+(x_1-x_3)^{2}\bigr\}
\]
である。$\vect{x}=(x_1,x_2,x_3)^{t}$ とし、$\zeta=\ee^{\ii2\pi/3}$、$U_{n\mu}=\zeta^{\mu n}/\sqrt{3}$（$n,\mu=1,2,3$）とする。
\begin{shoumon}
  \item $V=\frac{1}{2}\vect{x}^{t}S\vect{x}$ となる実対称行列 $S$ を求めよ。
  \item $U^{\dagger}SU$ を計算して対角行列になることを示し、その対角成分 $\lambda_1,\lambda_2,\lambda_3$ を求めよ。
  \item 運動方程式
\[
m\dv[2]{x_n}{t}=-\pdv{V}{x_n}
\]
を、$\hat{\vect{x}}=U^{\dagger}\vect{x}$ を用いて書き換えよ。各 $\hat{x}_{\mu}$ が互いに独立な単振動をすることを示し、その角振動数 $\omega_{\mu}$ を求めよ。
  \item 初期条件 $\vect{x}(0)=(d,0,0)^{t}$、$\dv{\vect{x}}{t}(0)=\vect{0}$ のもとで $x_1(t),x_2(t),x_3(t)$ を求めよ。さらに $x_1+x_2+x_3$ の時間変化を求め、その意味を述べよ。
\end{shoumon}
\notekiroku{2026-10-03 D ◎△××}
\begin{kotaeutsushi}{2026-10-03}
\toi{1}{◎}{}
\begin{align*}
  V&=\frac{\kappa}{2}(x_1^{2}+x_2^{2}+x_3^{2})\\
  &\qquad+\frac{k}{2}\bigl(2x_1^{2}+2x_2^{2}+2x_3^{2}\\
  &\qquad\kakimachigai{-x_1x_2-x_2x_3-x_3x_1}{-2(\cdots)}\bigr)
\end{align*}
\[
  S=\begin{pmatrix}\kappa+2k&-k&-k\\-k&\kappa+2k&-k\\-k&-k&\kappa+2k\end{pmatrix}
\]
\toi{2}{△}{計算}
（ページの上に固有方程式のメモ）
\[ (\kappa+2k-\lambda)^{3}-2k^{3}-3k^{2}(\kappa+2k+\lambda) \]
\[
  U=\frac{1}{\sqrt{3}}\begin{pmatrix}\ee^{\ii\frac{2\pi}{3}}&\ee^{\ii\frac{4\pi}{3}}&\ee^{\ii\frac{6\pi}{3}}\\[2pt]\ee^{\ii\frac{4\pi}{3}}&\ee^{\ii\frac{8\pi}{3}}&\ee^{\ii\frac{12\pi}{3}}\\[2pt]\ee^{\ii\frac{6\pi}{3}}&\ee^{\ii\frac{12\pi}{3}}&\ee^{\ii\frac{18\pi}{3}}\end{pmatrix}
\]
\begin{align*}
  U^{\dagger}SU&=\frac{1}{3}\begin{pmatrix}\zeta^{-1}&\zeta^{-2}&\zeta^{-3}\\\zeta^{-2}&\zeta^{-4}&\zeta^{-6}\\\zeta^{-3}&\zeta^{-6}&\zeta^{-9}\end{pmatrix}\\
  &\qquad\times\begin{pmatrix}\kappa+2k&-k&-k\\-k&\kappa+2k&-k\\-k&-k&\kappa+2k\end{pmatrix}\\
  &\qquad\times\begin{pmatrix}\zeta&\zeta^{2}&\zeta^{3}\\\zeta^{2}&\zeta^{4}&\zeta^{6}\\\zeta^{3}&\zeta^{6}&\zeta^{9}\end{pmatrix}\\
  &=\frac{1}{3}\begin{pmatrix}\zeta^{\fumei}(\kappa+2k)&\cdots\end{pmatrix}\nuke{\lambda_1,\lambda_2,\lambda_3}
\end{align*}
\naze{立式は正しいが，対角成分を出す前に止まっている。$S$ の各行は対角に $\kappa+2k$，両隣に $-k$ なので，$U_{n\pm1,\mu}=\zeta^{\pm\mu}U_{n\mu}$ を使うと $(SU)_{n\mu}=[\kappa+2k-k(\zeta^{\mu}+\zeta^{-\mu})]U_{n\mu}$ とまとまり，$\lambda_{\mu}=\kappa+2k-2k\cos(2\pi\mu/3)$，すなわち $\lambda_1=\lambda_2=\kappa+3k$，$\lambda_3=\kappa$ が出る。行列の積を全部書く必要はない。}
\toi{3}{×}{未着手}
\[
  \dv[2]{\hat{x}_1}{t}=-\frac{\kappa+3k}{m}\hat{x}_1,\qquad x=
\]
（図：3個の質点を三角形に並べ，ばねでつないだ図）
\naze{1つの式だけで止まっている。$\partial V/\partial x_n=(S\vect{x})_n$ より $m\,\dd^{2}\vect{x}/\dd t^{2}=-S\vect{x}$，$UU^{\dagger}=I$ を挟んで左から $U^{\dagger}$ をかけると $m\,\dd^{2}\hat{x}_{\mu}/\dd t^{2}=-\lambda_{\mu}\hat{x}_{\mu}$ となり，各モードが独立に単振動する。$\omega_1=\omega_2=\sqrt{(\kappa+3k)/m}$，$\omega_3=\sqrt{\kappa/m}$。}
\toi{4}{×}{未着手}
\mikai
\end{kotaeutsushi}
\begin{kizuki}
\item 環状の系の $U^{\dagger}SU$ は，行列の積を書き下さずに「$S$ の第 $n$ 行は対角 $\kappa+2k$，両隣 $-k$」と「$U_{n\pm1,\mu}=\zeta^{\pm\mu}U_{n\mu}$」から $(SU)_{n\mu}=\lambda_{\mu}U_{n\mu}$ を出す。$U^{\dagger}U=I$ と合わせれば $U^{\dagger}SU$ は対角で，対角成分が $\lambda_{\mu}$ になる。固有方程式の3次式を解くのは遠回りである。
\item 基準座標への書き換えは「$m\ddot{\vect{x}}=-S\vect{x}$ に左から $U^{\dagger}$ をかけ，$\vect{x}=U\hat{\vect{x}}$ を入れる」の1行で，全モードの式と $\omega_{\mu}=\sqrt{\lambda_{\mu}/m}$ が一度に出る。1つのモードだけ書いて終えない。
\end{kizuki}
\begin{tadashii}
\houshin 3質点のリング状ばね系のポテンシャルを2次形式で書き、ばね行列 $S$ を $N=3$ の離散フーリエ級数展開の行列 $U$ で対角化する。基準座標の運動方程式を解き、与えられた初期条件の解を元の座標に戻す。使う定理：2次形式 $V=\frac{1}{2}\vect{x}^{t}S\vect{x}$ と基準座標 $\hat{\vect{x}}=U^{\dagger}\vect{x}$。$m\dv[2]{\hat{x}_{\mu}}{t}=-\lambda_{\mu}\hat{x}_{\mu}$、$\omega_{\mu}=\sqrt{\lambda_{\mu}/m}$。$\zeta^{3}=1$、$\zeta+\zeta^{2}=-1$、$\zeta^{\mu}+\zeta^{-\mu}=2\cos\frac{2\pi\mu}{3}$。 % 出典: p.65--66
\keisan
\shoumonkai{1} 展開すると $(x_2-x_1)^{2}+(x_3-x_2)^{2}+(x_1-x_3)^{2}=2(x_1^{2}+x_2^{2}+x_3^{2})-2(x_1x_2+x_2x_3+x_3x_1)$ なので
\[
V=\frac{1}{2}\Bigl[(\kappa+2k)(x_1^{2}+x_2^{2}+x_3^{2})-2k(x_1x_2+x_2x_3+x_3x_1)\Bigr],\qquad
S=\begin{pmatrix}\kappa+2k&-k&-k\\-k&\kappa+2k&-k\\-k&-k&\kappa+2k\end{pmatrix}.
\]
\shoumonkai{2} $S$ の第 $n$ 行は、対角に $\kappa+2k$、残りの2つ（$n\pm1$、添字は3を法とする）に $-k$ をもつ。$\zeta^{3\mu}=1$ より $U_{n\pm3,\mu}=U_{n\mu}$、$U_{n\pm1,\mu}=\zeta^{\pm\mu}U_{n\mu}$ なので
\[
(SU)_{n\mu}=(\kappa+2k)U_{n\mu}-k\bigl(U_{n+1,\mu}+U_{n-1,\mu}\bigr)=\bigl[\kappa+2k-k(\zeta^{\mu}+\zeta^{-\mu})\bigr]U_{n\mu}=\lambda_{\mu}U_{n\mu},
\]
$\lambda_{\mu}=\kappa+2k-2k\cos\frac{2\pi\mu}{3}$。$U$ はユニタリーなので $(U^{\dagger}SU)_{\nu\mu}=\lambda_{\mu}(U^{\dagger}U)_{\nu\mu}=\lambda_{\mu}\delta_{\nu\mu}$ で対角。
$\cos(2\pi/3)=\cos(4\pi/3)=-1/2$、$\cos2\pi=1$ より $\lambda_1=\lambda_2=\kappa+3k$、$\lambda_3=\kappa$。
\shoumonkai{3} $\pdv{V}{x_n}=(S\vect{x})_n$ なので $m\dv[2]{\vect{x}}{t}=-S\vect{x}$。$UU^{\dagger}=I$ を挟み、左から $U^{\dagger}$ をかけると
\[
m\dv[2]{\hat{\vect{x}}}{t}=-(U^{\dagger}SU)(U^{\dagger}\vect{x})=-\mathrm{diag}(\lambda_1,\lambda_2,\lambda_3)\,\hat{\vect{x}},\qquad m\dv[2]{\hat{x}_{\mu}}{t}=-\lambda_{\mu}\hat{x}_{\mu}.
\]
各 $\hat{x}_{\mu}$ は他と結合しない単振動の方程式に従い、$\omega_{\mu}=\sqrt{\lambda_{\mu}/m}$。すなわち $\omega_1=\omega_2=\sqrt{(\kappa+3k)/m}\equiv\Omega$、$\omega_3=\sqrt{\kappa/m}\equiv\omega_0$。
\shoumonkai{4} 初速度は $\vect{0}$ なので
\[
\dv{\hat{\vect{x}}}{t}(0)=U^{\dagger}\dv{\vect{x}}{t}(0)=\vect{0}
\]
となり、$\hat{x}_{\mu}(t)=\hat{x}_{\mu}(0)\cos\omega_{\mu}t$。初期値は
\[
\hat{x}_{\mu}(0)=\sum_n U_{n\mu}^{*}x_n(0)=U_{1\mu}^{*}d=\frac{\zeta^{-\mu}}{\sqrt{3}}d .
\]
元の座標に戻すと
\[
x_n(t)=\sum_{\mu=1}^{3}U_{n\mu}\hat{x}_{\mu}(t)=\frac{d}{3}\sum_{\mu=1}^{3}\zeta^{\mu(n-1)}\cos\omega_{\mu}t .
\]
$n=1$: $\zeta^{0}=1$ より $x_1=\frac{d}{3}(2\cos\Omega t+\cos\omega_0t)$。
$n=2$: $\zeta^{\mu}$ は $\mu=1,2,3$ で $\zeta,\zeta^{2},1$。$\zeta+\zeta^{2}=-1$ より $x_2=\frac{d}{3}(\cos\omega_0t-\cos\Omega t)$。
$n=3$: $\zeta^{2\mu}$ は $\zeta^{2},\zeta,1$ で同様に $x_3=\frac{d}{3}(\cos\omega_0t-\cos\Omega t)$。
（確認: $t=0$ で $x_1=d$、$x_2=x_3=0$。）
和は $x_1+x_2+x_3=d\cos\omega_0t$。これは $\mu=3$（$U_{n3}=1/\sqrt{3}$ で全質点が同じ変位）のモードだけが和に寄与することを表す。質点どうしをつなぐばね $k$ の力は作用・反作用で打ち消し合うので、全体の平均の変位 $\frac{1}{3}(x_1+x_2+x_3)$ は釣り合いの位置へのばね $\kappa$ だけで決まる角振動数 $\omega_0=\sqrt{\kappa/m}$ で振動する。
\kotae{\shoumonkai{1} $S=\begin{pmatrix}\kappa+2k&-k&-k\\-k&\kappa+2k&-k\\-k&-k&\kappa+2k\end{pmatrix}$。\shoumonkai{2} $U^{\dagger}SU=\mathrm{diag}(\kappa+3k,\,\kappa+3k,\,\kappa)$。\shoumonkai{3} $m\dv[2]{\hat{x}_{\mu}}{t}=-\lambda_{\mu}\hat{x}_{\mu}$、$\omega_1=\omega_2=\sqrt{(\kappa+3k)/m}=\Omega$、$\omega_3=\sqrt{\kappa/m}=\omega_0$。\shoumonkai{4} $x_1=\frac{d}{3}(\cos\omega_0t+2\cos\Omega t)$、$x_2=x_3=\frac{d}{3}(\cos\omega_0t-\cos\Omega t)$、$x_1+x_2+x_3=d\cos\omega_0t$（全体の平均の変位は $\kappa$ のばねだけで振動する）。}
\kenzan $t=0$ で $x_1=d$、$x_2=x_3=0$ と初期条件を再現すること、$x_1+x_2+x_3=d\cos\omega_0t$ が $\kappa$ のばねだけで決まる振動になることを確かめた。また $k=0$ では $\Omega=\omega_0$ となり、$x_1=d\cos\omega_0t$、$x_2=x_3=0$ と、質点がそれぞれ独立に単振動する場合に一致する。$\lambda_{\mu}$ はばね定数と同じ次元で、$\omega_{\mu}=\sqrt{\lambda_{\mu}/m}$ は角振動数の次元をもつ。
\end{tadashii}
\end{mondai}

% toubun: kiroku\toubun\数理物理1\063_2026-10-03.tex
\mondaibangou{5}
\begin{mondai}{063-basis2}{計算}{2次元複素ベクトル空間の2組の正規直交基底}
2次元の複素ベクトル空間で Dirac の記法を用いる。$|1\rangle,|2\rangle$ は正規直交完備な基底（$\langle n|m\rangle=\delta_{nm}$）とする。新しいベクトルを
\[
|+\rangle=\frac{1}{\sqrt{2}}\bigl(|1\rangle+\ii|2\rangle\bigr),\qquad |-\rangle=\frac{1}{\sqrt{2}}\bigl(|1\rangle-\ii|2\rangle\bigr)
\]
で定める。
\begin{shoumon}
  \item $\langle+|+\rangle=\langle-|-\rangle=1$、$\langle+|-\rangle=0$ を示せ。
  \item 行列 $U_{n\mu}=\langle n|\mu\rangle$（$n=1,2$、$\mu=+,-$）を書き下し、$U^{\dagger}U=I$ を示せ。
  \item $|+\rangle\langle+|+|-\rangle\langle-|$ の行列要素 $\langle n|\bigl(|+\rangle\langle+|+|-\rangle\langle-|\bigr)|m\rangle$ を計算し、これが単位行列になること（完備性）を示せ。
  \item $|f\rangle=2|1\rangle+\ii|2\rangle$ について $\hat{f}_{\pm}=\langle\pm|f\rangle$ を求め、$|f\rangle=|+\rangle\hat{f}_{+}+|-\rangle\hat{f}_{-}$ と $\langle f|f\rangle=\abs{\hat{f}_{+}}^{2}+\abs{\hat{f}_{-}}^{2}$ を確かめよ。
\end{shoumon}
\notekiroku{2026-10-03 S ◎◎◎◎}
\begin{kotaeutsushi}{2026-10-03}
\toi{1}{◎}{}
\begin{align*}
  \langle+|+\rangle&=\frac{1}{\sqrt{2}}\bigl(\langle1|-\ii\langle2|\bigr)\cdot\frac{1}{\sqrt{2}}\bigl(|1\rangle+\ii|2\rangle\bigr)\\
  &=\frac{1}{2}\bigl(\langle1|1\rangle+\langle2|2\rangle\bigr)=1\\
  \langle-|-\rangle&=\frac{1}{\sqrt{2}}\bigl(\langle1|+\ii\langle2|\bigr)\cdot\frac{1}{\sqrt{2}}\bigl(|1\rangle-\ii|2\rangle\bigr)\\
  &=\frac{1}{2}\bigl(\langle1|1\rangle+\langle2|2\rangle\bigr)=1\\
  \langle+|-\rangle&=\frac{1}{\sqrt{2}}\bigl(\langle1|-\ii\langle2|\bigr)\cdot\frac{1}{\sqrt{2}}\bigl(|1\rangle-\ii|2\rangle\bigr)\\
  &=\frac{1}{2}\bigl(\langle1|1\rangle-\ii\langle2|1\rangle-\ii\langle1|2\rangle-\langle2|2\rangle\bigr)\\
  &=0
\end{align*}
\toi{2}{◎}{}
\begin{align*}
  \langle1|+\rangle&=\frac{1}{\sqrt{2}}\bigl(\langle1|1\rangle+\ii\langle1|2\rangle\bigr)=\frac{1}{\sqrt{2}}\\
  \langle1|-\rangle&=\frac{1}{\sqrt{2}}\bigl(\langle1|1\rangle-\ii\langle1|2\rangle\bigr)=\frac{1}{\sqrt{2}}\\
  \langle2|+\rangle&=\frac{1}{\sqrt{2}}\bigl(\langle2|1\rangle+\ii\langle2|2\rangle\bigr)=\frac{1}{\sqrt{2}}\ii\\
  \langle2|-\rangle&=\kakimachigai{\frac{1}{\sqrt{3}}}{\frac{1}{\sqrt{2}}}\kakimachigai{\,}{(}\langle2|1\rangle-\ii\langle2|2\rangle\bigr)=-\frac{1}{\sqrt{2}}\ii
\end{align*}
\[
  U_{n\mu}=\frac{1}{\sqrt{2}}\begin{pmatrix}1&1\\\ii&-\ii\end{pmatrix}
\]
\begin{align*}
  U^{\dagger}U&=\frac{1}{2}\begin{pmatrix}1&-\ii\\1&\ii\end{pmatrix}\begin{pmatrix}1&1\\\ii&-\ii\end{pmatrix}\\
  &=\frac{1}{2}\begin{pmatrix}2&0\\0&2\end{pmatrix}=I
\end{align*}
\toi{3}{◎}{}
\begin{align*}
  |+\rangle\langle+|&=\frac{1}{\sqrt{2}}\bigl(|1\rangle+\ii|2\rangle\bigr)\cdot\frac{1}{\sqrt{2}}\bigl(\langle1|-\ii\langle2|\bigr)\\
  &=\frac{1}{2}\bigl(|1\rangle\langle1|+|2\rangle\langle2|\\
  &\qquad-\ii|1\rangle\langle2|+\ii|2\rangle\langle1|\bigr)\\
  |-\rangle\langle-|&=\frac{1}{\sqrt{2}}\bigl(|1\rangle-\ii|2\rangle\bigr)\frac{1}{\sqrt{2}}\bigl(\langle1|+\ii\langle2|\bigr)\\
  &=\frac{1}{2}\bigl(|1\rangle\langle1|+|2\rangle\langle2|\\
  &\qquad-\ii|2\rangle\langle1|+\ii|1\rangle\langle2|\bigr)
\end{align*}
\[
  |+\rangle\langle+|+|-\rangle\langle-|=|1\rangle\langle1|+|2\rangle\langle2|
\]
\begin{align*}
  &\langle1|\bigl(|1\rangle\langle1|+|2\rangle\langle2|\bigr)|1\rangle\\
  &\qquad=\langle1|1\rangle\langle1|1\rangle+\langle1|2\rangle\langle2|1\rangle=1\\
  &\langle1|\bigl(|1\rangle\langle1|+|2\rangle\langle2|\bigr)|2\rangle\\
  &\qquad=\langle1|1\rangle\langle1|2\rangle+\langle1|2\rangle\langle2|2\rangle=0\\
  &\langle2|\bigl(|1\rangle\langle1|+|2\rangle\langle2|\bigr)|1\rangle\\
  &\qquad=\langle2|1\rangle\langle1|1\rangle+\langle2|2\rangle\langle2|1\rangle=0\\
  &\langle2|\bigl(|1\rangle\langle1|+|2\rangle\langle2|\bigr)|2\rangle\\
  &\qquad=\langle2|1\rangle\langle1|2\rangle+\langle2|2\rangle\langle2|2\rangle=1
\end{align*}
\toi{4}{◎}{}
\begin{align*}
  \hat{f}_{+}&=\langle+|f\rangle=2\langle+|1\rangle+\ii\langle+|2\rangle\\
  &=\sqrt{2}\bigl(\langle1|1\rangle-\ii\langle2|1\rangle\bigr)\\
  &\qquad+\frac{1}{\sqrt{2}}\ii\bigl(\langle1|2\rangle-\ii\langle2|2\rangle\bigr)\\
  &=\sqrt{2}+\frac{1}{\sqrt{2}}=\frac{3}{2}\sqrt{2}\\
  \hat{f}_{-}&=\langle-|f\rangle=2\langle-|1\rangle+\ii\langle-|2\rangle\\
  &=\sqrt{2}\bigl(\langle1|1\rangle+\ii\langle2|1\rangle\bigr)\\
  &\qquad+\frac{\ii}{\sqrt{2}}\bigl(\langle1|2\rangle+\ii\langle2|2\rangle\bigr)\\
  &=\sqrt{2}-\frac{1}{\sqrt{2}}=\frac{1}{2}\sqrt{2}
\end{align*}
\begin{align*}
  |f\rangle&=\Bigl(\sqrt{2}+\frac{1}{\sqrt{2}}\Bigr)\cdot\frac{1}{\sqrt{2}}\bigl(|1\rangle+\ii|2\rangle\bigr)\\
  &\qquad+\Bigl(\sqrt{2}-\frac{1}{\sqrt{2}}\Bigr)\cdot\frac{1}{\sqrt{2}}\bigl(|1\rangle-\ii|2\rangle\bigr)\\
  &=\frac{3}{2}\bigl(|1\rangle+\ii|2\rangle\bigr)+\frac{1}{2}\bigl(|1\rangle-\ii|2\rangle\bigr)\\
  &=2|1\rangle+\ii|2\rangle
\end{align*}
\begin{align*}
  \langle f|f\rangle&=\bigl(2\langle1|-\ii\langle2|\bigr)\bigl(2|1\rangle+\ii|2\rangle\bigr)\\
  &=4\langle1|1\rangle+2\ii\langle1|2\rangle-2\ii\langle2|1\rangle+\langle2|2\rangle=5
\end{align*}
\begin{align*}
  (\hat{f}_{+})^{2}+(\hat{f}_{-})^{2}&=\Bigl(\sqrt{2}+\frac{1}{\sqrt{2}}\Bigr)^{2}+\Bigl(\sqrt{2}-\frac{1}{\sqrt{2}}\Bigr)^{2}\\
  &=\Bigl(\frac{3}{2}\sqrt{2}\Bigr)^{2}+\Bigl(\frac{1}{2}\sqrt{2}\Bigr)^{2}=\frac{9}{2}+\frac{1}{2}=5
\end{align*}
\end{kotaeutsushi}
\end{mondai}

\noteowari
\end{document}
